\section{Introduction}
\label{sec:introduction}

Quantum computing has progressed from theoretical models of computation
to experimentally realized, programmable hardware, yet the currently available
noisy intermediate-scale quantum (NISQ) devices remain limited in qubit
count, connectivity, and coherence time~\citep{rehman2025evolutionary}.
This gap between the algorithms quantum computing ultimately promises and
the hardware available today has created a specific, practical need:
classical heuristics capable of tuning, compiling, and optimizing quantum
circuits and quantum-algorithm parameters well enough that near-term
devices can run them at all. Evolutionary algorithms (EAs) are a natural
fit for this role -- population-based, derivative-free, and already used
to design and compress circuits directly, as this survey's own case study
(Section~\ref{sec:case-study}) illustrates.

A second, independent thread of research asks the opposite question: can
ideas \emph{from} quantum mechanics -- superposition, interference,
entanglement -- improve classical evolutionary algorithms themselves, on
classical hardware, whether or not the resulting solutions ever touch a
real qubit? This is the domain of \emph{quantum-inspired evolutionary
algorithms} (QIEAs), founded on Han and Kim's Q-bit chromosome and
Q-gate rotation operator~\citep{han2002qea} and since applied across
combinatorial and continuous optimization
alike~\citep{rehman2025evolutionary,vivek2025fss}. Following
\citet{rehman2025evolutionary}, we adopt the terminology that separates
\emph{quantum-inspired} evolutionary algorithms (QiEA: classical
algorithms inspired by quantum phenomena) from \emph{quantum} evolutionary
algorithms (QEA: algorithms designed for and executed on physical quantum
hardware). This survey is scoped to the former.

Whether classical or quantum-inspired, evolutionary algorithms are above
all optimization methods -- yet the optimization problems they are
applied to seldom reduce to a single objective. Most real engineering and
scientific problems are not single-objective: they require trading
off several competing criteria (e.g.\ fidelity
against circuit depth, or cost against coverage) rather than optimizing
one metric in isolation. Multi-objective optimization (MOO) has its own
mature toolbox of classical evolutionary methods --
NSGA-II~\citep{deb2002nsga2}, NSGA-III~\citep{debjain2014nsga3},
SPEA2~\citep{zitzler2001spea2}, MOEA/D~\citep{zhangli2007moead}, and
SMS-EMOA~\citep{beume2007smsemoa} chief among them -- and QIEA research
has followed MOO into this territory since at least
\citet{kim2006qmea}'s quantum-inspired multi-objective knapsack algorithm.
Yet the literature applying QIEA specifically to multi-objective problems
remains comparatively sparse and scattered across
domains~\citep{olvera2023continuous,guzel2022qnsga2,kumar2026eahqnsga2},
compared to the much larger body of QIEA work on single- or
bi-objective problems such as feature subset
selection~\citep{zhang2011qieasurvey,vivek2025fss}. Part of the reason lies in the
representation itself: the Q-bit chromosome encodes each individual as
a vector of qubits that is observed into a binary string, a natural fit
for binary and combinatorial problems~\citep{han2002qea} but not for
problems whose solutions are more naturally encoded as permutations or
real-valued vectors, where an additional decoding step is required. As
a result, QIEAs do not always match classical EAs that operate on such
encodings directly: \citet{olvera2023continuous}, adapting two
originally combinatorial QIEA-MOO algorithms to continuous search
spaces, found that floating-point NSGA-II and SMPSO came closer to the
true Pareto front owing to their precision advantage, and that NSGA-II
converged faster overall, even though the quantum-inspired variants
were statistically competitive -- and that the quantum-inspired
variants' performance degraded as the number of decision variables,
and with it the number of Q-bits, increased. No survey, to our knowledge, systematically
taxonomizes and critically compares QIEA-for-MOO methods on their own
terms. This is the gap this survey addresses.

Beyond the mechanics of any one algorithm, the appeal of cheaper,
shallower, more efficient optimization has a stake that outlives any
individual benchmark result: computing's energy footprint. Large-scale
classical compute -- from generative-AI training runs to combinatorial
optimization on classical hardware -- carries real energy and
sustainability costs~\citep{sai2026generative}, and reducing the resources
a given optimization problem requires (fewer generations, cheaper fitness
evaluations, shallower resulting circuits) is not merely a performance
metric but part of a broader case for energy-efficient, sustainable
computing~\citep{banciu2026aiot}. Optimization-driven resource management
already delivers this kind of saving concretely outside quantum computing
-- reinforcement-learning-based energy management in community
microgrids, for instance, has been shown to meaningfully cut both cost
and carbon emissions relative to rule-based
control~\citep{moga2025rlmicrogrid} -- and the same underlying argument,
that a cheaper optimization process is itself a sustainability gain,
applies to this survey's subject matter. This is one of the underlying
reasons
QIEA-for-MOO is worth taking seriously as its own subfield rather than as
a footnote to single-objective QIEA or to classical MOO.

To ground the survey's more abstract taxonomy and critical-analysis
claims in a concrete example, Section~\ref{sec:case-study} presents a
case study built on this project's own quantum-circuit-optimization
pipeline, which partitions a target circuit into qubit blocks, evolves a
cheaper per-block circuit via a pluggable multi-objective optimizer, and
reassembles the result. That pipeline is itself a direct extension of
\citet{ghlib2026scalable}'s scalable, block-partitioned NSGA-II framework
for quantum circuit optimization, adding NSGA-III and SMS-EMOA as
alternative per-block optimizers. The case study is presented as an
illustrative worked example, not as this survey's primary evidence; full
empirical validation of that pipeline is the subject of a separate,
forthcoming empirical results paper.

\subsection{Scope and Inclusion Criteria}
\label{sec:introduction:scope}

This survey covers literature at the intersection of quantum-inspired
evolutionary computation and multi-objective optimization, identified
through a combination of targeted literature search (arXiv, IEEE Xplore,
Springer, ACM Digital Library, MDPI) and citation-chasing from the papers
found, spanning publications from 1985 (the earliest Pareto-based MOEA,
VEGA~\citep{schaffer1985vega}) to 2026. This targeted selection is
complemented by a bibliometric sweep of the QIEA literature as a whole,
2{,}250 papers from 1996--2026 retrieved from OpenAlex, which also
contributes its ten most-cited papers to the survey
(Section~\ref{sec:bibliometric}). We include:
\begin{itemize}
  \item Papers proposing or surveying quantum-inspired evolutionary
    algorithms (in the QiEA sense above) applied to multi-objective or
    many-objective optimization problems, on continuous or combinatorial
    search spaces.
  \item Classical (non-quantum-inspired) multi-objective evolutionary
    algorithms that QIEA-MOO methods are built on, hybridized with, or
    benchmarked against -- including the full Pareto-dominance lineage
    (VEGA through SPEA2 and beyond), indicator-based methods, and
    decomposition-based methods -- since these form the baseline
    vocabulary the taxonomy in Section~\ref{sec:taxonomy} depends on.
  \item Single-objective QIEAs, as a comparison baseline: the foundational
    representation and operator
    papers~\citep{narayanan1996qiga,han2002qea,han2004hepsilon} and
    broad single-objective QIEA
    surveys~\citep{zhang2011qieasurvey,vivek2025fss}. These are included
    not as QIEA-MOO methods in their own right but to show what changes,
    and what stays the same, when a QIEA moves from one objective to
    several (Sections~\ref{sec:taxonomy:objectives}
    and~\ref{sec:critical-analysis:so-to-mo}).
  \item Papers on (multi-objective) genetic algorithms for quantum
    circuit optimization specifically, since these motivate and inform
    the case study in Section~\ref{sec:case-study}.
  \item A small number of papers establishing the broader motivation for
    energy-efficient and sustainable computing, cited only in this
    introduction and the conclusion, and not treated as QIEA/MOO-technical
    literature.
\end{itemize}
We exclude quantum evolutionary algorithms designed to run on physical
quantum hardware (QEA in the strict sense), except where a paper
explicitly discusses the QiEA/QEA distinction. Only English-language
sources are included.

The multi-objective methods that the taxonomy
(Section~\ref{sec:taxonomy}) and feature comparison
(Section~\ref{sec:feature-comparison}) examine in depth are chosen by a
stated rule rather than by hand. They are the ten most-cited papers in
the bibliometric corpus of Section~\ref{sec:bibliometric} that are
tagged multi-objective and whose own optimization method is a QIEA --
skipping papers that use a QIEA only as a comparison baseline, and
papers without an English full text --
\citep{deng2020gate,li2007hqga,kim2011mqeaps,kim2006qmea,deng2024moqead,hussain2022moqiga,mousavi2019uav,wang2014powertrain,zhang2012moqeavrp,lu2013apmqea},
plus three recent methods found by the targeted search that a
citation-based rule would miss because they have had little time to be
cited~\citep{guzel2022qnsga2,olvera2023continuous,kumar2026eahqnsga2}.
This gives thirteen methods. For most of them, the comparison rests on
what the paper's abstract states; Section~\ref{sec:feature-comparison}
marks what could not be confirmed.

\subsection{Research Questions and Contributions}
\label{sec:introduction:contributions}

This survey is organized around four research questions:
\begin{description}
  \item[RQ1:] How can quantum-inspired evolutionary algorithms for
    multi-objective optimization be systematically classified?
  \item[RQ2:] How do the surveyed methods compare in representation,
    operators, selection mechanism, complexity, scalability, and
    reported performance?
  \item[RQ3:] What are the main limitations of the QIEA-for-MOO
    literature, both per method category and across the field as a
    whole?
  \item[RQ4:] How do these limitations manifest in a concrete
    multi-objective evolutionary pipeline applied to a quantum problem
    domain?
\end{description}
In answering them, this survey makes the following contributions, one
per research question, preceded by a bibliometric overview of the field:
\begin{itemize}
  \item A bibliometric analysis of 2{,}250 QIEA papers from 1996--2026:
    the field's publication trend, the share and recent growth of
    multi-objective work, its keyword clusters, and its ten most-cited
    papers (Section~\ref{sec:bibliometric}).
  \item \textbf{(RQ1)} A taxonomy of QIEA-for-MOO methods along five axes --
    representation scheme, quantum-inspired operator type,
    selection/diversity mechanism, hybridization with classical MOEAs,
    and application domain (Section~\ref{sec:taxonomy}).
  \item \textbf{(RQ2)} A systematic feature comparison of the surveyed methods against
    that taxonomy, covering representation, operator type, selection
    mechanism, and application domain alongside complexity, scalability,
    convergence properties, and reported performance
    (Section~\ref{sec:feature-comparison}).
  \item \textbf{(RQ3)} A critical analysis of per-category weaknesses and cross-cutting
    open problems in the field -- reproducibility gaps, weak
    benchmarking practices, and scalability claims that are rarely
    empirically stress-tested (Section~\ref{sec:critical-analysis}).
  \item \textbf{(RQ4)} A worked case study that grounds several of these abstract claims
    in a concrete block-partitioned quantum-circuit-optimization
    pipeline, itself an extension of prior published
    work~\citep{ghlib2026scalable} (Section~\ref{sec:case-study}).
\end{itemize}

\subsection{Paper Organization}
\label{sec:introduction:organization}

Section~\ref{sec:bibliometric} gives a bibliometric overview of the QIEA
literature. Section~\ref{sec:background} reviews the quantum-computing and
multi-objective-optimization background needed to follow the rest of the
survey, including the classical MOEA baselines QIEA methods are compared
against. Section~\ref{sec:taxonomy} introduces the classification axes
used throughout. Section~\ref{sec:feature-comparison} applies that
taxonomy systematically across the surveyed literature.
Section~\ref{sec:critical-analysis} discusses limitations and open
problems, per taxonomy category and across the field as a whole.
Section~\ref{sec:case-study} presents the illustrative case study.
Section~\ref{sec:open-challenges} synthesizes open challenges and future
directions, and Section~\ref{sec:conclusion} concludes.
